\documentclass[../main]{subfiles}
\usetikzlibrary{shapes,backgrounds}
\begin{document}
\chapter{Tarjetas Kanban y pizarras}
\img{images/02/kanban.jpg}{Tablero Kanban.}

\info{
Kanban}{
Se trata de un método visual de gestión de proyectos que permite a los equipos visualizar sus flujos y la carga de trabajo. En un tablero Kanban, el trabajo se muestra en forma de tablero organizado por columnas. Tradicionalmente, cada columna representa una etapa del trabajo. 
}

\section{Introducción}

En el desarrollo de proyectos, es común utilizar herramientas que permitan organizar las tareas de manera efectiva. Una de estas herramientas es la pizarra o tarjeta Kanban, que permite visualizar el estado de las tareas en un proyecto y mejorar la gestión del flujo de trabajo.

\section{¿Qué es una Pizarra o Tarjeta Kanban?}

La pizarra o tarjeta Kanban es una herramienta visual para la gestión de proyectos que se originó en la industria manufacturera japonesa. Consiste en un panel o pizarra en la que se colocan tarjetas o notas adhesivas que representan las tareas que se deben realizar en un proyecto.

\section{Ventajas de las Pizarras o Tarjetas Kanban}

Entre las ventajas que ofrecen las pizarras o tarjetas Kanban, se encuentran:

\begin{itemize}
    \item Visualización del estado de las tareas en tiempo real.
    \item Identificación de cuellos de botella y áreas de mejora en el flujo de trabajo.
    \item Facilidad para priorizar tareas y asignar responsabilidades.
    \item Mejora en la comunicación y colaboración del equipo.
\end{itemize}

\section{Ejemplo de Uso de una Pizarra o Tarjeta Kanban}

A continuación, se muestra un ejemplo de uso de una pizarra o tarjeta Kanban para la gestión de un proyecto de desarrollo de software:

En este ejemplo, las tarjetas se dividen en cuatro columnas: \textit{Backlog}, \textit{To Do}, \textit{Doing} y \textit{Done}. Cada tarjeta representa una tarea que debe ser realizada.

\actividad{Cuestionario de Comprensión}

\begin{enumerate}
    \item ¿Qué es Kanban y cómo funciona?
\item ¿Cuáles son los beneficios de utilizar Kanban en los procesos de trabajo?
\item ¿Cómo se implementa Kanban en diferentes contextos, como el desarrollo de software, la producción de manufactura, la gestión de proyectos, entre otros?
\end{enumerate}

\end{document}